\documentclass[11pt]{article}

% --- standard, widely-available packages (compiles with plain pdflatex) ---
\usepackage[utf8]{inputenc}
\usepackage[T1]{fontenc}
\usepackage{amsmath,amssymb}
\usepackage{booktabs}
\usepackage{array}
\usepackage[margin=1in]{geometry}
\usepackage{xcolor}
\usepackage[hidelinks]{hyperref}
\usepackage{enumitem}

\newcommand{\sig}{\operatorname{sig}}
\newcommand{\colour}{\chi}
\newcommand{\XXH}{\textsc{XXH3-64}}

\title{\bfseries Anchoring a Coordinate-Free Somatic Variant Graph:\\
Reference-Based Localization and Geometry-Driven\\ Reclassification of Reference-Free Structural-Variant Calls}

\author{svp project \\ \small (methodology note)}
\date{\today}

\begin{document}
\maketitle

\begin{abstract}
\noindent
\textbf{svp} calls somatic structural variants (SVs) directly from a two-colour
locally-consistent-parse (LCP) core graph built from normal and tumour long reads, without
assembly, alignment, or a reference genome. Its calls are therefore \emph{coordinate-free}:
each is a set of graph cores identified by hash, carrying no genomic position. We describe a
final, explicitly separated pipeline stage that \emph{anchors} these calls onto a reference
coordinate system so they can be emitted as a VCF, and that uses the recovered geometry to
correct systematic mis-typing produced by the coordinate-free caller. Anchoring maps each
call's reference-present flanking cores to unique genomic positions, brackets the breakpoint,
and then applies three signal-based operations: suppression of single-nucleotide ``mis-fires''
in the breakend caller, deduplication of the several redundant calls a real SV induces, and
geometry-driven reclassification of large deletions and insertions that the graph caller routes
through its breakend path. On a $\sim\!30\times$ whole-genome simulation these operations raise
precision from $4.3\%$ to $90.7\%$ and deletion recall from $82\%$ to $90\%$ at the
$\geq 50$\,bp scale (\textsc{Truvari} matching), with no parameter fitted to the truth set. We
further show, by extracting truth cores directly from the mutated genome, that the residual
recall gap is dominated by \emph{placement} (repeat-ambiguous flanks), not by detection: the
missed variants are present in the graph and largely called, but cannot be uniquely localized.
\end{abstract}

\section{Introduction}

Reference-free variant calling avoids the biases of mapping reads to a single linear reference,
at the cost of producing calls that live in a graph rather than on a coordinate axis. The
\textbf{svp} caller builds a two-colour LCP-core graph from normal and tumour high-fidelity (HiFi)
reads and emits somatic SV candidates as sets of graph cores. Downstream consumers, however, and
every standard SV benchmark, expect a VCF: chromosome, position, type, and length. The gap between
a coordinate-free call and a VCF record is the subject of this note.

We make three points. First, that gap is bridgeable using only the graph's own
reference-present cores as anchors, with no realignment of reads (Section~\ref{sec:anchor}).
Second, the recovered coordinates are not merely a reporting convenience: the \emph{geometry} of
a call's anchors---their spacing, and the colour of the reference cores between them---is a strong
signal that the coordinate-free caller lacks, and it repairs systematic type errors that the
graph alone cannot (Section~\ref{sec:refine}). Third, when a reference-free caller is evaluated,
the recall ceiling must be attributed carefully: we separate \emph{detection} (is the somatic
sequence in the graph?) from \emph{calling} (was a call emitted?) from \emph{placement} (could the
call be localized?), and find the residual loss is placement-bound (Section~\ref{sec:results}).

A methodological caveat pervades the evaluation and is reported honestly: truth ``somatic cores''
must be extracted from the mutated genome, not reconstructed by splicing an inserted sequence
between reference flanks, or insertion recall is silently undercounted (Section~\ref{sec:truth}).

\section{Background: the coordinate-free caller}
\label{sec:bg}

We summarize only what the anchoring stage consumes; the graph construction is prior work in the
pipeline.

\paragraph{Two-colour LCP-core graph.}
A locally-consistent parse partitions a sequence into variable-length \emph{cores} whose
boundaries depend only on a local window, so the same locus yields the same core in any context.
svp runs at level $\ell = 5$ (median core $\approx 150$\,bp; adjacent level-5 cores overlap
$\approx 85$\,bp). A node's identity is $\sig(c) = \XXH(c)$, a 64-bit hash of the core's bytes.
Reads from both samples are parsed on both strands, counted through a per-colour Bloom cascade
with a multiplicity floor, and survivors become nodes. Each node $v$ carries a colour
\begin{equation}
\colour(v) \;=\;
\begin{cases}
3 & \text{present in normal and tumour (shared),}\\
1 & \text{normal-only,}\\
2 & \text{tumour-only (candidate somatic).}
\end{cases}
\end{equation}
and per-base normal/tumour coverage. Edges join adjacent cores. The graph is coordinate-free:
node payload is core \emph{sequence}, never genome position.

\paragraph{Somatic calls.}
A somatic allele must involve a colour-2 core. The caller emits three kinds of call, each a
node set:
\begin{itemize}[nosep]
\item \textbf{BUBBLE} --- a two-arm bubble with a tumour arm; typed SNP/INS/DEL from the signed
spliced arm-length difference.
\item \textbf{BND} --- a tumour-only run bridging two shared anchors at near-full depth (a novel
adjacency); intended for DUP/INV/large-indel/translocation.
\item \textbf{TIP} --- a tumour-only dead-end chain (insertion-end candidate).
\end{itemize}
By design the caller is high-recall: it emits \emph{every} structure carrying a tumour-only allele
and defers false-positive control downstream. On our benchmark it emits $294{,}740$ calls
($262{,}951$ bubbles, $19{,}676$ breakends, $12{,}113$ tips) from a graph of $K = 86{,}519{,}562$
nodes. Each call also records the \XXH{} signatures of its alt-allele cores.

\section{Reference anchoring}
\label{sec:anchor}

Anchoring is an explicit, opt-in stage, \texttt{svp anchor <dir> <ref.fasta>}. It is deliberately
\emph{not} part of the reference-free \texttt{run} pipeline: detection stays reference-free, and a
reference is consulted only to name coordinates for reporting. The reference FASTA is the sole
external input; the mutated/tumour assembly is never used as a caller input.

\subsection{Anchor cores}
A somatic call is built from colour-2 (tumour-only) cores, which are by definition absent from the
reference and therefore unlocalizable on their own. Their graph neighbours, however, of colour
$1$ or $3$ are present in the normal genome and hence in the reference. For a call with alt-node
set $A$ we take as its \emph{anchor set}
\begin{equation}
\mathcal{N}(A) \;=\; \{\, w \in \mathcal{G} : \colour(w) \in \{1,3\},\ w \notin A,\
\exists\, a\in A\ \text{with}\ (a,w)\ \text{or}\ (w,a)\in E \,\}.
\end{equation}

\subsection{Reference core-position index}
We stream the reference once, parsing each contig on both strands at level $\ell$, and record, for
every needed anchor signature, its forward-genome position. Reverse-strand cores are mapped back
to forward coordinates. A signature seen at more than one distinct position is
\emph{repeat-ambiguous} and dropped, yielding a partial map
\begin{equation}
\pi:\ \sig \;\rightharpoonup\; (\text{contig}, \text{position}),
\end{equation}
that is defined only on cores with a unique genomic location. On the benchmark, $582{,}109$
distinct anchor signatures are required; $527{,}048$ receive a unique position and $36{,}107$ are
dropped as repeat-ambiguous.

\subsection{Breakpoint bracketing}
For each call we look up $\pi$ on its anchor set, keep the consensus contig (the one carrying the
most located anchors), and bracket the variant between the leftmost and rightmost anchor. A subtle
but important correction: an anchor's \emph{start} coordinate lies $\approx$ one core width
($\sim\!150$\,bp) upstream of the true breakpoint, because the breakpoint sits just past the left
flank core. We therefore place the left coordinate at the left anchor's \emph{end}
($\text{start} + \text{core length}$), removing a systematic $-170$\,bp offset (measured median)
that would otherwise defeat overlap-based matching for small SVs. A call all of whose anchors are
repeat-ambiguous cannot be placed; $20{,}186$ calls ($7.4\%$) are left unanchored.

\section{Signal-based refinement of anchored calls}
\label{sec:refine}

Coordinates unlock three operations the coordinate-free caller cannot perform. Let a breakend
call have anchor span $s$ (reference distance between its flanking anchors) and tumour junction
length $j$ (the spliced length of its tumour-only run).

\subsection{Suppression of single-nucleotide mis-fires (\textsc{LocalSub})}
\label{sec:localsub}
A single-nucleotide variant perturbs the graph locally; its tumour-only cores can bridge two
\emph{consecutive} reference anchors, which the breakend path then reports as a novel adjacency.
Such a ``breakend'' does not skip or rearrange the reference---its junction \emph{reproduces} the
spanned reference:
\begin{equation}
\text{mis-fire} \;\iff\; j \;\geq\; s/2 .
\end{equation}
The two regimes separate cleanly (Table~\ref{tab:ratio}): true DUP/INV breakends have a short
junction relative to their span (median $j/s = 0.02$), whereas mis-fires reproduce the span
(median $j/s = 1.04$). Flagged calls are marked \texttt{FILTER=LocalSub}; on the benchmark this
removes $7{,}666$ spurious breakends.

\begin{table}[t]
\centering
\begin{tabular}{lrrr}
\toprule
& median $j/s$ & $j/s \ge 0.7$ (local) & $j/s < 0.3$ (real skip) \\
\midrule
true DUP/INV breakends & $0.02$ & $6\%$ & $84\%$ \\
false breakends        & $1.04$ & $95\%$ & $4\%$ \\
\bottomrule
\end{tabular}
\caption{Junction-to-span ratio separates single-nucleotide mis-fires from true novel
adjacencies.}
\label{tab:ratio}
\end{table}

\subsection{Locus deduplication}
A single real SV induces several calls: a bubble \emph{and} a breakend, each doubled by the
forward/reverse-complement parse. Under one-to-one benchmark matching, one becomes a true positive
and the rest false positives. We collapse calls of the same type whose anchored positions fall
within a locus window, keeping the highest-support representative. This removes $43{,}767$
redundant records ($15.9\%$) and, as shown below, is the single largest precision gain.

\subsection{Geometry-driven reclassification}
Large deletions and insertions do not form clean bubbles (their reference arm branches through
repeats), so the caller routes them through the breakend path and mis-types them. The anchor
geometry recovers the correct type.

\paragraph{Deletion.}
A somatic deletion removes reference sequence present in the normal genome. Hence the reference
cores spanning $[\,\text{lo},\text{hi}\,]$ between the two anchors are predominantly
\emph{normal-only} (colour $1$) in the graph, whereas a duplication's span is shared (colour $3$)
and an artefact's is mixed. This test is strand-agnostic. For a two-anchor breakend with
$50 \le s \le 50{,}000$ and $j < s/2$ we fetch the spanned reference, parse it, and retype to
\textbf{DEL} of size $s$ when
\begin{equation}
\frac{\#\{\text{spanned cores with } \colour = 1\}}{\#\{\text{spanned cores in graph}\}} \;\ge\; 0.6 .
\end{equation}

\paragraph{Insertion.}
A somatic insertion places novel sequence between two \emph{adjacent} reference positions: its two
flanking anchors are essentially coincident ($s < 50$) while its junction is long
($j \ge 50$). We retype such calls to \textbf{INS} of size $\approx j - s$ \emph{before} the
\textsc{LocalSub} test, which would otherwise absorb them (a long junction trivially satisfies
$j \ge s/2$).

\section{Evaluation protocol}
\label{sec:eval}

\paragraph{Data.}
A human reference is mutated by insilicoSV to a truth genome and truth VCF ($125{,}830$ variants);
HiFi reads are simulated by Badread (normal from the reference, tumour from the truth genome) at
$\sim\!30\times$ whole-genome depth; svp is run at multiplicity floor $3$. The
$\geq 50$\,bp benchmark set contains $485$ SVs ($240$ DEL, $78$ INS, $167$ DUP$/$INV).

\paragraph{Matching.}
We benchmark with \textsc{Truvari} (sequence comparison disabled, as our symbolic alleles carry no
inserted sequence; breakend distance $1000$\,bp; \texttt{PASS}-only). Because the caller does not
sub-type DUP versus INV, we relabel the truth DUP/INV records to the breakend type so that
matching is type-consistent; deletion and insertion types are unchanged and \emph{not} matched
across type. No leniency (no type-ignoring) is used.

\paragraph{Detection accounting.}
\label{sec:truth}
To attribute the recall ceiling we build, per truth variant, the set of its \emph{out-only} cores
(present in the tumour genome, absent from the reference) and test membership in (i) the read
graph and (ii) the emitted call set. We stress a construction pitfall: reconstructing an
insertion's cores as [reference flank]$+$[inserted sequence]$+$[reference flank] yields junction
cores that do not match the tumour haplotype when other variants perturb the flanks, so the
insertion appears to produce \emph{zero} somatic cores. Extracting the cores from the mutated
genome directly removes the artefact: all sampled ``zero-core'' insertions in fact produce
$6$--$10$ novel cores, all present in the read graph.

\section{Results}
\label{sec:results}

\subsection{Ablation}
Table~\ref{tab:ablation} adds each operation in turn. Anchoring alone gives high recall but
catastrophic precision, because every real SV is over-called and the breakend path mis-fires on
the dense SNPs. \textsc{LocalSub} and deduplication together lift precision from $4.3\%$ to
$84.3\%$; geometry-driven retyping recovers deletions (and their precision) and slightly raises
insertion handling, reaching $90.7\%$ precision at $78.4\%$ recall.

\begin{table}[t]
\centering
\begin{tabular}{lrrr}
\toprule
configuration & recall & precision & $F_1$ \\
\midrule
anchored calls (no refinement)            & $77.1\%$ & $4.3\%$  & $0.08$ \\
\; $+$ \textsc{LocalSub} mis-fire filter  & $75.5\%$ & $37.9\%$ & $0.51$ \\
\; $+$ locus deduplication                & $74.2\%$ & $84.3\%$ & $0.79$ \\
\; $+$ BND$\to$DEL retyping                & $78.4\%$ & $90.9\%$ & $0.84$ \\
\; $+$ BND$\to$INS retyping (final)        & $78.4\%$ & $90.7\%$ & $0.84$ \\
\bottomrule
\end{tabular}
\caption{Contribution of each anchoring operation at the $\geq 50$\,bp scale (\textsc{Truvari};
$485$ truth SVs). No parameter is fitted to the truth set.}
\label{tab:ablation}
\end{table}

\subsection{Per-class performance}
Table~\ref{tab:perclass} gives the final per-class recall. Deletion recall reaches $90\%$;
DUP$/$INV and insertion remain lower for reasons attributed below.

\begin{table}[t]
\centering
\begin{tabular}{lrr}
\toprule
class & recalled / truth & recall \\
\midrule
DEL          & $216/240$ & $90.0\%$ \\
INS          & $56/78$   & $71.8\%$ \\
DUP$/$INV    & $108/167$ & $64.7\%$ \\
\midrule
all $\geq 50$\,bp & $380/485$ & $78.4\%$ \\
\bottomrule
\end{tabular}
\caption{Final per-class recall (precision $90.7\%$ overall, $380$ TP / $39$ FP).}
\label{tab:perclass}
\end{table}

\subsection{Detection versus placement}
The residual recall gap is not a detection failure. Of the $24$ missed deletions, $18$ are
\emph{called} (an out-only core appears in an emitted call) but cannot be placed because their
flanking anchors are repeat-ambiguous; only $6$ are true calling misses. Of the $22$ missed
insertions, $21$ likewise have no anchored call near their locus despite being detected in the
graph. Deletion cores are detected at $97.2\%$ ($280/288$) and insertion cores, measured against
tumour-genome-derived truth, at essentially $100\%$. The ceiling is therefore \emph{placement}:
variants whose flanks fall in repeats have no unique coordinate for any reference-based anchoring
to assign.

\section{Discussion and limitations}

\paragraph{Coordinates repair type, not just report it.}
The most transferable observation is that reference geometry is an independent signal channel. The
same anchor spacing that yields a VCF position also distinguishes a deletion (normal-only span)
from a duplication (shared span) and a point insertion (coincident anchors, long junction) from a
substitution (reproduced span). A coordinate-free caller cannot access these, so a modest,
principled anchoring stage recovers accuracy that no amount of graph-only reasoning would.

\paragraph{The placement wall.}
Because anchoring relies on uniquely-locatable flanking cores, variants embedded in repeats are
lost even when fully detected and called. Raising deletion and insertion recall past $\sim\!90\%$
would require multi-position anchoring---retaining a repeat-ambiguous anchor's several candidate
positions and placing a call where its anchors mutually agree---at some risk to precision. This is
the natural next step and is left as future work.

\paragraph{Sub-typing.}
DUP versus INV is not resolved; both are reported as breakends. Distinguishing them requires the
oriented order of the two anchors (forward for deletion, reversed for tandem duplication, opposite
strand for inversion), which in turn requires tracking anchor strand through the dual
forward/reverse-complement parse.

\paragraph{Honest accounting.}
Aggregate recall on this data is dominated by SNPs ($94\%$ of truth variants) and is not
informative about SVs; per-class, size-stratified figures are the only meaningful ones. Truth
cores for insertions must be taken from the mutated genome, or insertion detection is
underreported---an artefact we encountered and corrected.

\section{Conclusion}

We localize the coordinate-free somatic calls of a reference-free SV caller onto a reference using
only their graph-present flanking cores, and show that the recovered geometry both produces a VCF
and repairs the caller's systematic mis-typing. On a whole-genome simulation the anchoring stage
lifts $\geq 50$\,bp precision from $4.3\%$ to $90.7\%$ and deletion recall to $90\%$ without any
truth-fitted parameter, and it localizes the remaining ceiling to variants in repeats that no
reference-based method can uniquely place. Detection is not the bottleneck; placement is.

\end{document}
